\documentclass[11pt,a4paper]{article}

\usepackage[margin=1in]{geometry}
\usepackage{amsmath,amssymb,amsthm}
\usepackage{graphicx}
\usepackage{booktabs}
\usepackage{microtype}
\usepackage{xcolor}
\usepackage[hidelinks,colorlinks=true,linkcolor=blue!50!black,citecolor=blue!40!black,urlcolor=blue!50!black]{hyperref}
\usepackage[numbers,sort&compress]{natbib}

\newtheorem{theorem}{Theorem}
\newtheorem{proposition}[theorem]{Proposition}
\newtheorem{lemma}[theorem]{Lemma}
\newtheorem{corollary}[theorem]{Corollary}
\theoremstyle{definition}
\newtheorem{definition}[theorem]{Definition}
\newtheorem{conditional}[theorem]{Conditional Proposition}
\newtheorem{classical}[theorem]{Classical Theorem}
\newtheorem{interpretation}[theorem]{Foundational Interpretation}
\theoremstyle{remark}
\newtheorem{remark}[theorem]{Remark}

\title{\textbf{Toward the Euler Bridge}\\[2pt]
\large From Closure-Irreducibles to the Classical $\zeta$ Product via $\zeta_{\mathbb{Q}(\sqrt{5})}(s)$\\[2pt]
\normalsize \emph{Closure-Irreducibles, Paper III of a paired-plus-bridge contribution.}\\[1pt]
\small Companions: Paper~I (\emph{The Critical Line as Observation Manifold}) and Paper~II (\emph{Classical Primes as Closure-Irreducibles}).\\[1pt]
\small (Open research note --- pre-peer-review draft)}

\author{%
  Lee Smart\\[4pt]
  \textit{Institute of Vibrational Field Dynamics}\\[2pt]
  \texttt{contact@vibrationalfielddynamics.org}\\[2pt]
  \texttt{@vfd\_org}%
}
\date{May 2026}

\begin{document}
\maketitle
\thispagestyle{empty}

\begin{abstract}
\noindent
\textbf{Scope.} Papers~I and~II~\cite{CriticalLinePaperI,ClosurePrimesPaperII}
placed classical integer primes within the closure-irreducibility
typology as level $L_4$, using the classical Dedekind splitting law in
$\mathbb{Z}[\varphi]$. The present paper asks the natural follow-on
question: \emph{when one assembles those $L_4$ closure-irreducibles into
their natural Euler product, does the classical Riemann $\zeta(s)$
appear?}
\textbf{Method.} Two ingredients, both classical and unconditional. (i)
By Paper~II, the closure-irreducibles at $L_4$ are exactly the prime
ideals of $\mathbb{Z}[\varphi]$. By the standard definition of Dedekind
$\zeta$, their Euler product is $\zeta_{\mathbb{Q}(\sqrt{5})}(s)$. (ii)
By classical algebraic number theory (Dedekind, Hecke, textbook for
over a century), this Dedekind $\zeta$ factors as
$\zeta_{\mathbb{Q}(\sqrt{5})}(s) = \zeta(s) \cdot L(s, \chi_5)$, where
$\chi_5$ is the Kronecker symbol mod $5$.
\textbf{Primary structural result.} The classical Riemann $\zeta(s)$
appears as the $\sigma$-symmetric Euler factor of the cascade's natural
$L_4$ $\zeta$. The Dirichlet $L$-function $L(s, \chi_5)$ is the
$\sigma$-antisymmetric Euler factor, encoding the cascade $\sigma$-pairing
at the character level.
\textbf{Status.} The two underlying mathematical statements are
unconditional and classical. Their interpretation as the Euler bridge from
cascade $L_4$ closure-irreducibles to classical $\zeta$ is a structural
reading, conditional on the named hypotheses of Paper~II
(H-SigmaRestrict, H-UniqueSigma, H-NaturalIrreducibles) plus one new
named hypothesis H-EulerBridge: that the natural $\zeta$ attached to the
$L_4$ closure-irreducibles is the Dedekind $\zeta_{\mathbb{Q}(\sqrt{5})}$
and not some other arithmetic invariant.
\textbf{What this provides.} An explicit Euler-product placement of
$\zeta(s)$ within the cascade $L_4$ infrastructure, sim-verified at
$s = 2, 3, 4$ to ten or more decimal digits, and a structural reading of
the $\sigma$-symmetric / $\sigma$-antisymmetric decomposition of
$\zeta_{\mathbb{Q}(\sqrt{5})}(s)$ as the trace of the $L_4$ Galois
involution.
\textbf{What this does not provide.} A proof or conditional reduction of
the classical Riemann Hypothesis; a derivation of $\zeta(s)$ from
cascade machinery (it is an extracted factor, not a constructed object);
a bridge from the cascade analytic object
$\hat{\zeta}(z)$~\cite{rhFormalSmart} to classical $\zeta(s)$ (still
blocked by the Hecke obstruction~\cite{rhFormalSmart}); a verification
that $\zeta(s)$ has the same zero locations as $L(s, \chi_5)$ (it does
not; they have distinct zero sets).
\textbf{Relation to classical results.} The Dedekind factorisation
$\zeta_K(s) = \zeta(s) \cdot L(s, \chi)$ for quadratic fields is
classical and unconditional. The contribution of the present paper is
its \emph{interpretation} as the Euler bridge from cascade $L_4$ to
classical $\zeta$, not its derivation.
\end{abstract}

\bigskip

\noindent\fbox{%
\parbox{0.96\linewidth}{%
\textbf{Plain-language hypothesis.}

\smallskip\noindent
Paper~II showed where the classical primes sit in the cascade
typology: they are level $L_4$, classified by the
$\mathbb{Q}(\sqrt{5})$ Galois symmetry $\sqrt{5} \mapsto -\sqrt{5}$
into split, inert, and ramified primes.

\smallskip\noindent
The present paper asks the natural next question: \emph{if you take
those $L_4$ closure-irreducibles and form their Euler product the way
you form any zeta function, does classical Riemann $\zeta(s)$ appear?}

\smallskip\noindent
The answer, modulo named structural hypotheses, is \textbf{yes, but as
a factor, not as the whole thing}:
\[
\underbrace{\zeta_{\mathbb{Q}(\sqrt{5})}(s)}_{\text{cascade's natural $L_4$ $\zeta$}}
\;=\;
\underbrace{\zeta(s)}_{\text{$\sigma$-symmetric factor}}
\;\cdot\;
\underbrace{L(s, \chi_5)}_{\text{$\sigma$-antisymmetric factor}}.
\]
This is classical Dedekind factorisation. What is new is the
\emph{interpretation} that classical $\zeta(s)$ is the
$\sigma$-symmetric trace inside the cascade's natural $L_4$ $\zeta$,
and $L(s, \chi_5)$ is the $\sigma$-antisymmetric trace.

\smallskip\noindent
This is the \emph{Euler bridge}: a structural placement of classical
$\zeta(s)$ inside the cascade infrastructure. It is \emph{not}:
\begin{itemize}\itemsep0pt
  \item A proof or conditional reduction of the Riemann Hypothesis.
  \item A derivation of $\zeta(s)$ from cascade alone (the factor is
        extracted, not constructed).
  \item A bridge from cascade's analytic $\hat{\zeta}(z)$ to classical
        $\zeta(s)$ --- that bridge is still blocked by the Hecke
        obstruction.
\end{itemize}
}}

\bigskip

\paragraph{Programme map.}
This paper is the third in the Closure-Irreducibles series. Paper~I
introduces the five-level typology of closure-irreducibles and the
critical-line / observation-manifold foundation. Paper~II develops
$L_4$ via the Dedekind splitting law in $\mathbb{Z}[\varphi]$ and
establishes the universal $\sigma$-pattern across $L_0$--$L_4$. The
present paper takes the $L_4$ closure-irreducibles of Paper~II and
forms their natural Euler product, identifying classical Riemann
$\zeta(s)$ as one of two factors.

\paragraph{Status taxonomy.}
\textbf{Theorem} = formal structural claim under stated assumptions.
\textbf{Classical Theorem} = result imported from classical number theory,
cited with reference (Dedekind, Hecke, textbook). \textbf{Conditional
Proposition} = bridge claim depending on a named hypothesis.
\textbf{Foundational Interpretation} = structural reading of classical
mathematical content within the closure framework, marked as such.

\paragraph{Non-claims (load-bearing).}
This paper does \textbf{not} claim that:
\begin{itemize}\itemsep0pt
  \item Classical $\zeta(s)$ is derived from cascade machinery (it is extracted as a factor; the construction of the factor depends on classical Dedekind theory).
  \item The Riemann Hypothesis is proved, addressed at zero-location level, or conditionally reduced (it is not; the bridge concerns Euler-product factorisation, not zero locations).
  \item Cascade's analytic object $\hat{\zeta}(z)$ equals $\zeta(s)$ (it does not; the Hecke obstruction blocks this~\cite{rhFormalSmart}).
  \item $\zeta(s)$ and $L(s, \chi_5)$ have the same zero set (they do not; only the GRH for both would place them on the same critical line, and neither is proved).
  \item The Kronecker symbol $\chi_5$ is a new mathematical object (it is classical, going back to Kronecker; what is new is its interpretation as the cascade's $\sigma$-antisymmetric trace).
\end{itemize}
The framework is one of \emph{interpretation through the closure lens},
not of redefinition or proof.

\section{Introduction: the Euler bridge question}\label{sec:intro}

Paper~I established the closure-irreducibility typology with five
demonstrated levels $L_0$--$L_4$, and identified the critical line
$\mathrm{Re}(s) = 1/2$ as the fixed locus of the Mellin involution under
the observation hypothesis H-Observe. Paper~II developed $L_4$ as the
classical integer primes, classified by the Galois conjugation
$\sigma_K : \sqrt{5} \mapsto -\sqrt{5}$ in $\mathbb{Q}(\sqrt{5})$ via the
classical Dedekind splitting law.

A natural follow-on question, raised but not addressed in either earlier
paper, is whether the $L_4$ closure-irreducibles \emph{assemble into a
$\zeta$ function}, and if so, how that $\zeta$ relates to the classical
Riemann $\zeta(s)$.

The question has a definite answer, drawing on two unconditional inputs:

\begin{enumerate}
\item By Paper~II, the closure-irreducibles at level $L_4$ are
      (under H-NaturalIrreducibles) the prime ideals of
      $\mathbb{Z}[\varphi]$. The natural Euler product over these primes
      is, by definition, the Dedekind $\zeta$-function of the field
      $\mathbb{Q}(\sqrt{5})$:
      \begin{equation}\label{eq:dedekind_def}
        \zeta_{\mathbb{Q}(\sqrt{5})}(s)
          = \prod_{\mathfrak{p} \subset \mathbb{Z}[\varphi]}
              \frac{1}{1 - N(\mathfrak{p})^{-s}},
      \end{equation}
      where the product runs over nonzero prime ideals and $N(\mathfrak{p})$
      is the absolute norm.
\item By classical algebraic number theory (Dedekind 1871; modern
      reference Neukirch~\cite{Neukirch1999}, Davenport~\cite{Davenport2000}),
      the Dedekind $\zeta$ of any quadratic field $\mathbb{Q}(\sqrt{d})$
      factors as
      \begin{equation}\label{eq:dedekind_factorisation}
        \zeta_{\mathbb{Q}(\sqrt{d})}(s)
          = \zeta(s) \cdot L(s, \chi_d),
      \end{equation}
      where $\chi_d$ is the quadratic Dirichlet character associated to
      the discriminant of the field. For $d = 5$, the discriminant is
      $5$ and $\chi_5$ is the Kronecker symbol $\left(\tfrac{5}{n}\right)$
      (equivalently, the Legendre symbol when $n$ is an odd prime
      $\neq 5$).
\end{enumerate}

Combining (\ref{eq:dedekind_def}) and (\ref{eq:dedekind_factorisation})
gives the Euler bridge:
\begin{equation}\label{eq:euler_bridge}
  \zeta_{L_4}(s)
    \;\equiv\;
    \zeta_{\mathbb{Q}(\sqrt{5})}(s)
    \;=\;
    \zeta(s) \cdot L(s, \chi_5).
\end{equation}

The classical Riemann $\zeta(s)$ appears as one of two Euler factors
of the cascade's natural $L_4$ $\zeta$. The companion factor
$L(s, \chi_5)$ is the Dirichlet $L$-function attached to the quadratic
character $\chi_5$.

This paper develops the structural content of (\ref{eq:euler_bridge}):
identifies $\zeta(s)$ as the $\sigma$-symmetric Euler factor and
$L(s, \chi_5)$ as the $\sigma$-antisymmetric Euler factor in the
cascade $L_4$ Galois decomposition; states explicitly the named
hypotheses on which the cascade interpretation depends; sim-verifies
the factorisation numerically at $s = 2, 3, 4$; and demarcates clearly
what the bridge does and does not provide.

\section{The cascade's natural $L_4$ $\zeta$}\label{sec:l4_zeta}

\subsection{Definition by Euler product over $L_4$ closure-irreducibles}

\begin{definition}[$\zeta_{L_4}$, cascade $L_4$ $\zeta$]\label{def:zeta_L4}
Let $\mathcal{P}_{L_4}$ denote the closure-irreducibles at level $L_4$
of the typology of Paper~I~\cite{CriticalLinePaperI}, as developed in
Paper~II~\cite{ClosurePrimesPaperII}. The \emph{cascade $L_4$
$\zeta$-function} is the formal Euler product
\begin{equation}\label{eq:zeta_L4_def}
  \zeta_{L_4}(s) \;:=\;
    \prod_{\mathfrak{q} \in \mathcal{P}_{L_4}}
      \frac{1}{1 - \mathrm{nm}(\mathfrak{q})^{-s}},
\end{equation}
where $\mathrm{nm}(\mathfrak{q})$ denotes the closure-norm of
$\mathfrak{q}$ (Paper~II, \S 3).
\end{definition}

This definition is purely formal at the level of the cascade. To
identify $\zeta_{L_4}(s)$ with a concrete analytic object, we must
identify $\mathcal{P}_{L_4}$ and the closure-norm with a concrete
classical structure. Paper~II provides exactly this identification.

\subsection{Identification with Dedekind $\zeta_{\mathbb{Q}(\sqrt{5})}$}

\begin{conditional}[$\zeta_{L_4}$ equals Dedekind $\zeta_{\mathbb{Q}(\sqrt{5})}$, under H-NaturalIrreducibles + H-EulerBridge]\label{prop:l4_dedekind}
\emph{Assume H-NaturalIrreducibles} (Paper~II, \S 5.3): the cascade
closure-irreducibles at $L_4$ are the prime ideals of
$\mathbb{Z}[\varphi]$, and the closure-norm is the absolute norm
$N : \mathfrak{p} \mapsto |\mathbb{Z}[\varphi]/\mathfrak{p}|$.

\emph{Assume H-EulerBridge}: the natural $\zeta$-function attached to
the $L_4$ closure-irreducibles is their formal Euler product
(\ref{eq:zeta_L4_def}), and not some other arithmetic invariant (e.g.\
a twisted, weighted, or geometric $\zeta$).

\smallskip Then
\begin{equation}
  \zeta_{L_4}(s) \;=\; \zeta_{\mathbb{Q}(\sqrt{5})}(s)
\end{equation}
in the sense of formal Euler products, with convergence for
$\mathrm{Re}(s) > 1$ and meromorphic continuation given by the
classical theory of Dedekind $\zeta$-functions
\cite{Neukirch1999,Marcus1977}.
\end{conditional}

\begin{proof}[Proof sketch]
Under H-NaturalIrreducibles, $\mathcal{P}_{L_4}$ is the set of prime
ideals of $\mathbb{Z}[\varphi]$, and the closure-norm coincides with the
absolute norm. Under H-EulerBridge, $\zeta_{L_4}(s)$ is the formal
Euler product (\ref{eq:zeta_L4_def}). The right-hand side of
(\ref{eq:zeta_L4_def}) is then literally the defining Euler product of
$\zeta_{\mathbb{Q}(\sqrt{5})}(s)$ (compare equation
(\ref{eq:dedekind_def})). The remaining content --- absolute convergence
for $\mathrm{Re}(s) > 1$, meromorphic continuation, functional equation
--- is classical Dedekind theory.
\end{proof}

\begin{remark}
The named hypothesis H-EulerBridge is a discipline hypothesis, not a
mathematical leap: it says only that the cascade's natural $\zeta$ at
$L_4$ is the Euler product over $L_4$ closure-irreducibles, rather than
some twist, weighting, or alternative construction. It is the
counterpart, at the $\zeta$-function level, of the more familiar
identification ``the natural $\zeta$ of a number field is its Dedekind
$\zeta$.'' We name it explicitly so that any future alternative cascade
$\zeta$ (for instance, a weighted Euler product reflecting closure-flow
multiplicities) can be addressed without ambiguity about which one the
present bridge concerns.
\end{remark}

\subsection{Status of $\zeta_{L_4}$ vs.\ $\hat{\zeta}$}

It is important to keep $\zeta_{L_4}(s)$ distinct from the cascade
analytic object $\hat{\zeta}(z)$ studied in \emph{rh-formal}
\cite{rhFormalSmart}. These are two different analytic functions:

\begin{itemize}
\item $\zeta_{L_4}(s)$ is an Euler product over the $L_4$
      closure-irreducibles (this paper). Under
      Conditional~Proposition~\ref{prop:l4_dedekind}, it equals
      $\zeta_{\mathbb{Q}(\sqrt{5})}(s)$.
\item $\hat{\zeta}(z)$ is the pentagonal-clock Mellin object with
      $K$-poles at $K \in \{0, 20, 52, 72\}$ with multiplicities
      $\{1, 1, 5, 5\}$ (\emph{rh-formal} Theorem $M1$).
\end{itemize}

These objects are structurally related (both arise from the cascade
$\sigma$-paired class infrastructure) but they are not the same
analytic function. In particular, the Hecke obstruction
(\emph{rh-formal} Theorem $M3$) shows that the multiplicities
$\{1,1,5,5\}$ of $\hat{\zeta}(z)$ do not match classical Hecke
multiplicities and so $\hat{\zeta}(z)$ cannot equal $\zeta(s)$ as
analytic objects. This obstruction does \emph{not} affect the present
bridge, because the present bridge concerns $\zeta_{L_4}(s)$, not
$\hat{\zeta}(z)$.

\section{Classical Dedekind factorisation for quadratic fields}\label{sec:dedekind_fact}

We now import the second ingredient of the bridge: the classical
factorisation of Dedekind $\zeta$ for a quadratic field.

\begin{classical}[Dedekind factorisation for quadratic fields, Dedekind 1871]\label{thm:dedekind_factorisation}
Let $K = \mathbb{Q}(\sqrt{d})$ with $d$ a squarefree integer, and let
$D_K$ denote the discriminant of $K$ (so $D_K = d$ if $d \equiv 1
\pmod 4$ and $D_K = 4d$ otherwise). Let $\chi_{D_K}$ denote the
Kronecker symbol $\left(\tfrac{D_K}{\cdot}\right)$, which is the unique
real primitive Dirichlet character of conductor $|D_K|$. Then
\begin{equation}
  \zeta_K(s) \;=\; \zeta(s) \cdot L(s, \chi_{D_K})
\end{equation}
as analytic functions on $\mathrm{Re}(s) > 1$, extending to a
meromorphic identity on $\mathbb{C}$ via the classical functional
equations of both factors.
\end{classical}

\begin{proof}[Reference]
This is textbook material; see Neukirch~\cite[Ch.~VII, \S5]{Neukirch1999},
Marcus~\cite[Ch.~7]{Marcus1977}, or
Davenport~\cite[Ch.~5]{Davenport2000}. The factorisation is a special
case of the general decomposition of $\zeta_K$ into Hecke $L$-functions
indexed by characters of $\mathrm{Gal}(K/\mathbb{Q})$; for $K$ quadratic,
$\mathrm{Gal}(K/\mathbb{Q}) = \{1, \sigma_K\}$ and the two characters are
the trivial character (yielding $\zeta(s)$) and the sign character
(yielding $L(s, \chi_{D_K})$).
\end{proof}

\subsection{Specialisation to $K = \mathbb{Q}(\sqrt{5})$}

For $d = 5$: $5 \equiv 1 \pmod 4$, hence $D_K = 5$ and $\chi_5$ is the
Kronecker symbol modulo $5$. Explicitly,
\[
  \chi_5(n) \;=\; \left(\tfrac{n}{5}\right)
  \;=\;
  \begin{cases}
    +1 & \text{if } n \equiv \pm 1 \pmod{5},\\
    -1 & \text{if } n \equiv \pm 2 \pmod{5},\\
    \phantom{+}0 & \text{if } n \equiv 0 \pmod{5}.
  \end{cases}
\]
The Dirichlet $L$-function is then
\[
  L(s, \chi_5) \;=\; \sum_{n=1}^{\infty} \frac{\chi_5(n)}{n^s}
    \;=\; \prod_{p \neq 5} \frac{1}{1 - \chi_5(p) p^{-s}}.
\]
Applying Classical Theorem~\ref{thm:dedekind_factorisation} at $d = 5$
gives
\begin{equation}\label{eq:factorisation_q5}
  \zeta_{\mathbb{Q}(\sqrt{5})}(s) \;=\; \zeta(s) \cdot L(s, \chi_5).
\end{equation}

\subsection{Note on the trivial character interpretation}

The factor $\zeta(s)$ in (\ref{eq:factorisation_q5}) is exactly the
classical Riemann $\zeta$-function. It is the Dirichlet $L$-function
attached to the trivial character $\chi_0$ in this decomposition:
$L(s, \chi_0) = \zeta(s)$. This is the standard interpretation and the
starting point of all classical work connecting $\zeta_K$ to $\zeta$;
we mention it explicitly because the cascade interpretation of the next
section makes use of it.

\section{The Euler bridge: $\zeta(s)$ as the $\sigma$-symmetric factor}\label{sec:euler_bridge}

We can now combine the cascade identification of \S\ref{sec:l4_zeta}
with the classical factorisation of \S\ref{sec:dedekind_fact} to state
the Euler bridge.

\subsection{Main bridge theorem}

\begin{conditional}[Euler bridge: $\zeta(s)$ inside cascade $L_4$, under H-NaturalIrreducibles + H-EulerBridge]\label{thm:euler_bridge}
Under the named hypotheses H-NaturalIrreducibles (Paper~II, \S 5.3) and
H-EulerBridge (Conditional Proposition~\ref{prop:l4_dedekind}), the
cascade $L_4$ $\zeta$-function factors as
\begin{equation}\label{eq:bridge_main}
  \zeta_{L_4}(s) \;=\; \zeta(s) \cdot L(s, \chi_5),
\end{equation}
where $\zeta(s)$ is the classical Riemann $\zeta$-function and
$L(s, \chi_5)$ is the Dirichlet $L$-function attached to the
Kronecker symbol $\chi_5$ modulo $5$. The classical Riemann $\zeta(s)$
therefore appears as one of two Euler factors of the cascade $L_4$
$\zeta$.
\end{conditional}

\begin{proof}
Combine Conditional Proposition~\ref{prop:l4_dedekind}
($\zeta_{L_4}(s) = \zeta_{\mathbb{Q}(\sqrt{5})}(s)$) with Classical
Theorem~\ref{thm:dedekind_factorisation} at $d = 5$
($\zeta_{\mathbb{Q}(\sqrt{5})}(s) = \zeta(s) \cdot L(s, \chi_5)$).
\end{proof}

\subsection{The $\sigma$-symmetric / $\sigma$-antisymmetric decomposition}

The cascade interpretation of the two factors on the right-hand side
of (\ref{eq:bridge_main}) is as follows. The Galois group
$\mathrm{Gal}(\mathbb{Q}(\sqrt{5})/\mathbb{Q}) = \{1, \sigma_K\}$ acts
on Dirichlet characters of $\mathbb{Z}[\varphi]$ via $\sigma_K$. There
are exactly two characters of this group: the trivial character
(invariant under $\sigma_K$, the \emph{$\sigma$-symmetric} character)
and the sign character (sent to its negation under $\sigma_K$, the
\emph{$\sigma$-antisymmetric} character).

In the Hecke decomposition of $\zeta_K(s)$,
\begin{equation}\label{eq:hecke_decomp}
  \zeta_K(s) \;=\; \prod_{\chi \in \widehat{\mathrm{Gal}(K/\mathbb{Q})}}
    L(s, \chi),
\end{equation}
each character contributes one $L$-factor. The trivial character
contributes $\zeta(s)$ and the sign character contributes
$L(s, \chi_5)$.

\begin{interpretation}[$\zeta(s)$ as $\sigma$-symmetric trace, $L(s, \chi_5)$ as $\sigma$-antisymmetric trace]\label{interp:sigma_decomp}
Under the cascade interpretation, the Hecke decomposition
(\ref{eq:hecke_decomp}) for $K = \mathbb{Q}(\sqrt{5})$ is read as the
$\sigma$-symmetric / $\sigma$-antisymmetric splitting of the cascade
$L_4$ $\zeta$ under the cascade involution $\sigma_K$ (identified with
the Galois conjugation via H-SigmaRestrict, Paper~II, \S 3.2). The
classical Riemann $\zeta(s)$ is the $\sigma$-symmetric trace; the
Dirichlet $L$-function $L(s, \chi_5)$ is the $\sigma$-antisymmetric
trace.

\smallskip\noindent
This is a foundational interpretation, not a new mathematical theorem.
The Hecke decomposition itself is classical (Hecke 1920s, textbook).
What the cascade interpretation adds is the structural reading of
``trivial character'' as ``$\sigma$-symmetric'' and ``sign character''
as ``$\sigma$-antisymmetric,'' consistent with the identifications
H-SigmaRestrict and H-UniqueSigma of Paper~II.
\end{interpretation}

\subsection{Compatibility with the $L_4$ splitting classes}

The Euler-product reading of (\ref{eq:bridge_main}) is consistent with
the $L_4$ classification of Paper~II at the level of local Euler
factors. For a rational prime $p$, the local Euler factor on the
left-hand side of (\ref{eq:bridge_main}) depends on the splitting
behaviour of $p$ in $\mathbb{Z}[\varphi]$:

\begin{itemize}
\item \textbf{Split} ($p \equiv 1, 4 \pmod 5$, $\sigma$-paired):
      $p = \mathfrak{p}_1 \mathfrak{p}_2$ with two prime ideals
      $\mathfrak{p}_1, \mathfrak{p}_2$ of norm $p$. The local factor is
      $(1 - p^{-s})^{-2}$.
\item \textbf{Inert} ($p \equiv 2, 3 \pmod 5$, $\sigma$-fixed):
      $p\mathbb{Z}[\varphi]$ remains prime, with norm $p^2$. The local
      factor is $(1 - p^{-2s})^{-1}$.
\item \textbf{Ramified} ($p = 5$, ramified anchor):
      $5\mathbb{Z}[\varphi] = \mathfrak{p}^2$ with $\mathfrak{p}$ of
      norm $5$. The local factor is $(1 - 5^{-s})^{-1}$.
\end{itemize}

On the right-hand side, the $p$-th local factors of $\zeta(s)$ and
$L(s, \chi_5)$ multiply to give the same local factor in each case:
\begin{itemize}
\item \textbf{Split}: $(1 - p^{-s})^{-1} \cdot (1 - p^{-s})^{-1}
      = (1 - p^{-s})^{-2}$, using $\chi_5(p) = +1$.
\item \textbf{Inert}: $(1 - p^{-s})^{-1} \cdot (1 + p^{-s})^{-1}
      = (1 - p^{-2s})^{-1}$, using $\chi_5(p) = -1$.
\item \textbf{Ramified}: $(1 - 5^{-s})^{-1} \cdot 1
      = (1 - 5^{-s})^{-1}$, using $\chi_5(5) = 0$.
\end{itemize}

The compatibility at the local-factor level is the local-global content
of the bridge: each $L_4$ $\sigma$-class contributes the matching
local Euler factor on both sides.

\section{The named hypotheses and what they buy}\label{sec:hypotheses}

The bridge claim (\ref{eq:bridge_main}) rests on three named hypotheses,
two imported from Paper~II and one new to the present paper. We state
each explicitly and identify the cost of its failure.

\subsection{H-NaturalIrreducibles (imported from Paper~II)}

\textbf{Statement.} The cascade closure-irreducibles at $L_4$ are the
prime ideals of $\mathbb{Z}[\varphi]$ (not, for instance, the rational
primes paired with their splitting data, or some other arithmetic
object).

\textbf{Cost of failure.} If $L_4$ closure-irreducibles are not exactly
$\mathbb{Z}[\varphi]$-primes, the cascade $L_4$ $\zeta$ is not the
Dedekind $\zeta_{\mathbb{Q}(\sqrt{5})}$, and the bridge would need to
re-target a different classical $\zeta$ (if any).

\subsection{H-EulerBridge (new to this paper)}

\textbf{Statement.} The natural $\zeta$-function attached to the $L_4$
closure-irreducibles is the formal Euler product (\ref{eq:zeta_L4_def})
with closure-norm = absolute norm, not a twisted, weighted, or
geometric variant.

\textbf{Cost of failure.} If the natural $L_4$ $\zeta$ is a weighted
Euler product (for instance, with closure-flow multiplicities at each
prime), the bridge would still produce an Euler product but it would
not equal the Dedekind $\zeta_{\mathbb{Q}(\sqrt{5})}$, and the
factorisation into $\zeta(s) \cdot L(s, \chi_5)$ would not directly
apply. A weighted bridge could still be analysed, but the present paper
treats only the unweighted case.

\subsection{H-SigmaRestrict (imported from Paper~II)}

\textbf{Statement.} The cascade spectral involution $\tau_{\mathrm{spec}}$,
restricted to coefficients in $\mathbb{Q}(\sqrt{5})$, induces the
Galois conjugation $\sigma_K : \sqrt{5} \mapsto -\sqrt{5}$.

\textbf{Cost of failure.} The Hecke decomposition
(\ref{eq:hecke_decomp}) is classical regardless of cascade
interpretation. If H-SigmaRestrict fails, the structural reading of
$\zeta(s)$ as ``$\sigma$-symmetric trace'' (Interpretation~\ref{interp:sigma_decomp})
becomes an unmotivated relabeling rather than a structural
identification. The mathematical content of
Conditional~Proposition~\ref{thm:euler_bridge} is unaffected, but the
\emph{cascade interpretation} of the bridge as the
$\sigma$-symmetric / $\sigma$-antisymmetric decomposition collapses to
the classical (uninterpreted) Hecke decomposition.

\subsection{Summary table}

\begin{table}[h]
\centering
\renewcommand{\arraystretch}{1.2}
\begin{tabular}{lll}
\toprule
\textbf{Hypothesis} & \textbf{Source} & \textbf{Buys} \\
\midrule
H-NaturalIrreducibles & Paper II \S 5.3 & $\mathcal{P}_{L_4} = \mathrm{Spec}(\mathbb{Z}[\varphi])$ primes \\
H-EulerBridge & This paper \S 2.2 & $\zeta_{L_4}(s) = \zeta_{\mathbb{Q}(\sqrt{5})}(s)$ \\
H-SigmaRestrict & Paper II \S 3.2 & $\sigma$-symmetric / antisymmetric reading \\
\bottomrule
\end{tabular}
\caption{Named hypotheses behind the Euler bridge (\ref{eq:bridge_main})
and the structural content each one buys. The factorisation
$\zeta_K(s) = \zeta(s) \cdot L(s, \chi_5)$ is unconditional classical
mathematics.}\label{tab:hypotheses}
\end{table}

\section{Numerical verification}\label{sec:numerics}

The factorisation (\ref{eq:bridge_main}) is verified numerically by
the script \texttt{repro/verify\_euler\_bridge.py} accompanying this
paper. The script computes both sides of (\ref{eq:bridge_main}) by
independent Euler products and compares them to high precision.

\subsection{Method}

For each test value $s \in \{2, 3, 4\}$:
\begin{enumerate}
\item Compute $\zeta_{\mathbb{Q}(\sqrt{5})}(s)$ via the
      $\mathbb{Z}[\varphi]$ Euler product (\ref{eq:dedekind_def}),
      using the $L_4$ classification of Paper~II to enumerate prime
      ideals up to a chosen norm cutoff $N \le N_{\max}$.
\item Independently compute $\zeta(s)$ via the rational Euler product
      $\prod_{p \le P_{\max}} (1 - p^{-s})^{-1}$ for the same value of
      $s$.
\item Independently compute $L(s, \chi_5)$ via the $\chi_5$-twisted
      rational Euler product
      $\prod_{p \le P_{\max}, p \neq 5} (1 - \chi_5(p) p^{-s})^{-1}$.
\item Compare $\zeta_{\mathbb{Q}(\sqrt{5})}(s)$ against
      $\zeta(s) \cdot L(s, \chi_5)$.
\end{enumerate}

\subsection{Expected outcome}

For each $s > 1$, the three Euler products converge absolutely and
their values are related by (\ref{eq:bridge_main}). At
$P_{\max} = 10^5$, $N_{\max} = 10^5$, agreement to ten or more decimal
digits is expected and is the threshold built into the sim script.

\subsection{Outcome (sim-verified)}

The script confirms agreement to the expected precision at all three
test values $s \in \{2, 3, 4\}$, as recorded in
\texttt{repro/output/euler\_bridge\_verification.json}. Detailed
numerics are deferred to \S\ref{sec:reproduction}.

\section{What this bridge does and does not provide}\label{sec:scope}

\subsection{What it does provide}

\begin{itemize}
\item An explicit Euler-product placement of classical $\zeta(s)$
      within the cascade $L_4$ infrastructure
      (Conditional~Proposition~\ref{thm:euler_bridge}).
\item A structural reading of the Dedekind factorisation
      $\zeta_K = \zeta \cdot L(s, \chi_5)$ as the
      $\sigma$-symmetric / $\sigma$-antisymmetric decomposition of the
      cascade $L_4$ $\zeta$
      (Interpretation~\ref{interp:sigma_decomp}).
\item Numerical verification of the factorisation at
      $s = 2, 3, 4$ to ten or more decimal digits
      (\S\ref{sec:numerics}, \S\ref{sec:reproduction}).
\item Local-factor compatibility with the $L_4$ split / inert /
      ramified classification of Paper~II at every rational prime $p$
      (\S\ref{sec:euler_bridge}).
\end{itemize}

\subsection{What it does not provide}

\begin{itemize}
\item \textbf{Not a proof of the Riemann Hypothesis.} The factorisation
      (\ref{eq:bridge_main}) is about Euler products, not about zero
      locations. It is consistent with $\zeta(s)$ having any zero
      distribution whatsoever in the critical strip. Knowing that
      $\zeta(s)$ is one factor of $\zeta_{\mathbb{Q}(\sqrt{5})}(s)$
      does not in itself locate the zeros of either.
\item \textbf{Not a conditional reduction of RH to GRH for $\chi_5$.}
      Such a reduction would require a structural statement constraining
      $\zeta(s)$ zeros from $L(s, \chi_5)$ zeros, or vice versa, beyond
      what the trivial implication ``zeros of the product = zeros of one
      factor or the other'' provides. The cascade interpretation does
      not provide such a constraint.
\item \textbf{Not a derivation of $\zeta(s)$ from cascade machinery.}
      The factor $\zeta(s)$ is extracted from $\zeta_{\mathbb{Q}(\sqrt{5})}(s)$
      via classical Dedekind factorisation. The cascade does not
      \emph{construct} $\zeta(s)$; it identifies $\zeta(s)$ as the
      $\sigma$-symmetric factor of an Euler product that includes
      classical $\zeta$ as one component.
\item \textbf{Not a bridge from $\hat{\zeta}(z)$ to $\zeta(s)$.} The
      pentagonal-clock $\hat{\zeta}(z)$ of \emph{rh-formal} remains a
      distinct analytic object, with multiplicities $\{1,1,5,5\}$ at its
      $K$-poles that do not match Hecke multiplicities
      \cite{rhFormalSmart}. The Hecke obstruction blocks the
      $\hat{\zeta} \leftrightarrow \zeta$ identification, and the
      present bridge does not engage with that obstruction.
\item \textbf{Not a claim about the zero set of $L(s, \chi_5)$.} The
      generalised Riemann Hypothesis for $L(s, \chi_5)$ remains open.
      The cascade interpretation of $L(s, \chi_5)$ as the
      $\sigma$-antisymmetric trace does not address GRH.
\end{itemize}

\subsection{The Hecke obstruction does not block the present bridge}

\begin{remark}[Why $M3$ does not apply]\label{rem:m3}
The Hecke obstruction of \emph{rh-formal} $M3$~\cite{rhFormalSmart} is
the statement that the pentagonal-clock multiplicities $\{1,1,5,5\}$ of
$\hat{\zeta}(z)$ at $K \in \{0, 20, 52, 72\}$ do not match the Hecke
multiplicity structure of any classical $L$-function, so $\hat{\zeta}$
cannot be identified with $\zeta$ (or with any classical $L$-function)
as an analytic object. The present bridge does not assert any such
identification. The cascade $L_4$ $\zeta$ defined here is
$\zeta_{\mathbb{Q}(\sqrt{5})}$, which has a perfectly classical Hecke
multiplicity structure. The Hecke obstruction applies to
$\hat{\zeta}(z)$ (which sits at the cascade level $L_3$ of the
typology, the pentagonal-clock generating function) and not to
$\zeta_{L_4}(s)$ (which is the present paper's object). The two are
related as separate cascade analytic objects living at adjacent levels
of the typology; the bridge is between $\zeta_{L_4}$ and classical
$\zeta$, not between $\hat{\zeta}$ and classical $\zeta$.
\end{remark}

\section{Relation to existing structural approaches}\label{sec:relation}

\subsection{Beyond endoscopy (Langlands)}

Langlands's beyond-endoscopy programme~\cite{LanglandsBeyond} proposes
to study the spectral side of trace formulas as a source for $L$-function
poles. The present bridge does not engage with the spectral / geometric
trace-formula side; it works entirely at the Euler-product level. The
two approaches are in principle compatible but not directly comparable.

\subsection{Adelic / noncommutative geometry (Connes)}

Connes's adelic / noncommutative-geometric reformulation of RH
\cite{ConnesAdelic} positions $\zeta(s)$ as a spectral trace on an
adelic space. Again, this engages zero locations directly via spectral
theory, whereas the present bridge engages Euler-product structure
only. The present interpretation of $\zeta(s)$ as $\sigma$-symmetric
trace at the Euler-product level is in principle compatible with a
spectral interpretation but is at a different level of the argument.

\subsection{Hilbert--P\'olya-style spectral conjectures}

Hilbert--P\'olya-style approaches conjecture an operator whose spectrum
is the imaginary part of the non-trivial zeros of $\zeta(s)$
\cite{KeatingSnaith,BerryKeating1999}. The present bridge says nothing
about such operators; it identifies $\zeta(s)$ as an Euler factor of a
larger Dedekind $\zeta$, without analysing zeros.

\subsection{The Selberg trace formula}

The Selberg trace formula~\cite{Selberg1956} provides a duality
between geometric (closed geodesic) data and spectral (Laplacian
eigenvalue) data on hyperbolic surfaces. The cascade typology shares
a structural feature with Selberg's setup --- the duality between
spectral irreducibles ($L_2$) and arithmetic irreducibles ($L_4$) ---
but the explicit identification at the trace-formula level is not
pursued here.

\subsection{Comparison summary}

The present bridge is structurally narrower than each of the above
approaches: it does not engage zero locations, spectral operators, or
trace formulas. Its contribution is the explicit placement of classical
$\zeta(s)$ as one of two Euler factors of the cascade's natural $L_4$
$\zeta$, with both factors interpretable structurally
($\sigma$-symmetric and $\sigma$-antisymmetric trace) under the named
hypotheses imported from Paper~II.

\section{Falsifiers and limitations}\label{sec:falsifiers}

The bridge of this paper can be falsified or rendered vacuous in the
following ways. We list them explicitly for the same reason as Papers~I
and~II: a paper that does not name its falsifiers does not earn the
right to be examined.

\begin{enumerate}
\item \textbf{Failure of H-NaturalIrreducibles.} If the cascade $L_4$
      closure-irreducibles can be shown to be a different arithmetic
      object than the $\mathbb{Z}[\varphi]$-primes (for example, a
      different ring's primes, or rational primes paired with cascade
      flow data), the identification $\zeta_{L_4} = \zeta_K$ fails, and
      with it the Euler bridge in the form stated here.
\item \textbf{Failure of H-EulerBridge.} If the cascade's natural $L_4$
      $\zeta$ is shown to be a weighted or twisted Euler product rather
      than the unweighted one in (\ref{eq:zeta_L4_def}), the bridge in
      this paper would re-target a twisted $L$-function rather than
      $\zeta_K$, and the structural reading would shift accordingly.
\item \textbf{Failure of H-SigmaRestrict.} If the cascade spectral
      involution cannot be identified with $\sigma_K$, the
      $\sigma$-symmetric / $\sigma$-antisymmetric interpretation of
      $\zeta(s)$ and $L(s, \chi_5)$ becomes a relabeling of trivial /
      sign characters, without structural cascade content. The
      mathematical identity (\ref{eq:bridge_main}) is unaffected.
\item \textbf{Numerical sim disagreement.} If the sim of
      \S\ref{sec:reproduction} shows the two sides of
      (\ref{eq:bridge_main}) disagreeing at any of the tested $s$ values
      beyond the truncation-error bound, the underlying classical
      Dedekind factorisation would be in dispute --- which would be
      surprising given the century-plus of classical work, but the sim
      check exists precisely to make this falsifier concrete and
      computable.
\end{enumerate}

\section{Reproduction}\label{sec:reproduction}

The classical Dedekind factorisation (\ref{eq:factorisation_q5}) and
hence the Euler bridge (\ref{eq:bridge_main}) are sim-verified by the
script \texttt{repro/verify\_euler\_bridge.py} accompanying this paper
(numpy only, deterministic, $\sim$300 lines). The script:
\begin{itemize}
\item Sieves rational primes $p \le P_{\max} = 10^5$ and classifies
      each as split / inert / ramified in $\mathbb{Z}[\varphi]$ (using
      Paper~II classification, $p \bmod 5$).
\item Computes $\zeta_{\mathbb{Q}(\sqrt{5})}(s)$ via the
      $\mathbb{Z}[\varphi]$-Euler product, enumerating prime ideals of
      $\mathbb{Z}[\varphi]$ by case:
      $(1 - p^{-s})^{-2}$ for split, $(1 - p^{-2s})^{-1}$ for inert,
      $(1 - 5^{-s})^{-1}$ for ramified.
\item Independently computes $\zeta(s)$ via the rational Euler product
      $\prod_{p \le P_{\max}} (1 - p^{-s})^{-1}$.
\item Independently computes $L(s, \chi_5)$ via the $\chi_5$-twisted
      Euler product
      $\prod_{p \neq 5, p \le P_{\max}} (1 - \chi_5(p) p^{-s})^{-1}$.
\item Compares $\zeta_{\mathbb{Q}(\sqrt{5})}(s)$ against
      $\zeta(s) \cdot L(s, \chi_5)$ at $s \in \{2, 3, 4\}$ to 10+
      decimal digits.
\end{itemize}

The script is deterministic and produces JSON output and one figure.
All Euler-product factorisation checks pass at $\ge 10$ decimal digits
of agreement at $s = 2$ (and progressively higher precision at
$s = 3, 4$). The check is purely classical-arithmetic; it uses no
cascade machinery. Its role is to make the classical factorisation
visibly false-falsifiable: a single numerical disagreement would either
break the factorisation or expose a bug in the sim.

\section{Conclusion}\label{sec:conclusion}

The closure-irreducibles series, taken together, places classical
Riemann $\zeta(s)$ within the cascade infrastructure as follows:

\begin{itemize}
\item \textbf{Paper I:} Five-level typology $L_0$--$L_4$ of
      closure-irreducibles; critical line $\mathrm{Re}(s) = 1/2$ as
      $\mathrm{Fix}(\tau_{\mathrm{crit}})$.
\item \textbf{Paper II:} $L_4$ = classical integer primes via the
      Dedekind splitting law in $\mathbb{Z}[\varphi]$; universal
      $\sigma$-pattern across $L_0$--$L_4$.
\item \textbf{Paper III (the present paper):} Euler product of $L_4$
      closure-irreducibles is $\zeta_{\mathbb{Q}(\sqrt{5})}(s)$;
      classical Dedekind factorisation gives $\zeta(s)$ as the
      $\sigma$-symmetric Euler factor.
\end{itemize}

The placement is structural, not derivational. Classical $\zeta(s)$
remains classical $\zeta(s)$, with its zero distribution still the
subject of the Riemann Hypothesis. What the closure framework adds is
\emph{location}: $\zeta(s)$ is the $\sigma$-symmetric Euler factor of
the cascade's natural $L_4$ $\zeta$, and $L(s, \chi_5)$ is its
$\sigma$-antisymmetric companion. The Hecke obstruction that blocks
identification of cascade's pentagonal-clock $\hat{\zeta}(z)$ with
$\zeta(s)$ remains where it is; it lives at a different level of the
typology and is not in conflict with the present bridge.

The honest scope of the trilogy:
\begin{itemize}
\item \emph{We have not proved RH.} The bridge here concerns
      Euler-product factorisation, not zero locations.
\item \emph{We have not derived $\zeta(s)$.} The factor $\zeta(s)$ is
      extracted from $\zeta_{\mathbb{Q}(\sqrt{5})}(s)$ via classical
      Dedekind factorisation; the cascade identifies its position, not
      its existence.
\item \emph{We have placed $\zeta(s)$.} Within the cascade typology,
      classical Riemann $\zeta$ is now structurally located: it is the
      $\sigma$-symmetric Euler factor of the natural $L_4$ $\zeta$.
\end{itemize}

The trilogy is, in this sense, a structural placement programme rather
than a Millennium-attack programme. We regard the placement itself ---
classical primes as $L_4$ closure-irreducibles, classical $\zeta$ as
their $\sigma$-symmetric Euler factor --- as the contribution.

\paragraph{Acknowledgements.}
This paper builds on classical algebraic number theory (Dedekind,
Hecke, Davenport, Neukirch, Marcus), on Papers~I and~II of the present
series, and on the broader Existence / Life / Closure programme. The
author thanks the classical number-theory community whose century of
work on Dedekind $\zeta$-functions and Dirichlet $L$-functions is the
foundation on which this structural reinterpretation rests.

\bibliographystyle{plainnat}
\bibliography{references}

\end{document}
